\subsection{In-Context Segmentation}\label{sec:fundamentals-icl}

Given a target image $I^t$ and a context set
$\{(I^c_k, L^c_k)\}_{k=1}^{K}$ of image/mask pairs for the same class,
in-context segmentation predicts a mask $\hat{L}^t$ for the target in a
single forward pass, with no gradient update and no class-specific
parameters. This sits between two more familiar paradigms: closed-set
supervised segmentation, which fixes a label set at training time and
must be retrained (or fine-tuned) to add a class; and interactive
segmentation, where a human supplies clicks or boxes on the target image
itself, per image, at inference time. In-context segmentation instead
supplies the ``prompt'' as a handful of \emph{other}, already-labeled
examples of the same class -- closer to how a radiologist might be shown
a few example cases of a rare finding and asked to find it in a new scan.
This is valuable in medical imaging specifically because annotation
requires domain expertise, new structures and pathologies are defined
faster than segmentation datasets can be built for them, and per-site
protocol differences limit how far a closed-set model trained elsewhere
transfers.
